\section{Approximation with controlled cone factors}\label{sec:conditioned-approximation}
In an exact lift, a proper cone of dimension $m$ supplies at most
$\max\{m-2,0\}$ mixed-curvature directions. The bound extends to
approximate lifts only under additional control of the factors. We give
two versions: a theorem for finitely many sampled directions, which
uses no derivatives, and a local theorem for smooth factorizations, in
which the mixed derivative of the slack is exact. The hypotheses concern
the selected factors in fixed Euclidean coordinates and are stronger
than approximation of the projected convex body.

\subsection{A finite sampling theorem}
Write $r=N-1\geq1$, fix $0<h<1$, and put
\[
 v_0=e_N,\qquad v_j=he_j+\sqrt{1-h^2}e_N\quad(1\leq j\leq r),
 \qquad \tau=1-\sqrt{1-h^2}.
\]
Suppose proper cones $K_i\subset E_i$, of dimensions $m_i$, admit
factors $A_i(v_j)\in K_i$, $B_i(v_j)\in K_i^*$ such that
\begin{equation}\label{eq:finite-slack-error}
 \left|\sum_i\ip{A_i(v_j)}{B_i(v_\ell)}
       -(1-\ip{v_j}{v_\ell})\right|\leq\epsilon
 \qquad(0\leq j,\ell\leq r).
\end{equation}
Let $a_i=A_i(v_0)$, $b_i=B_i(v_0)$, and let $X_i,Y_i:\R^r\to E_i$
have columns $A_i(v_j)-a_i$, $B_i(v_j)-b_i$, respectively.
All matrix norms below are operator norms. Define
\begin{equation}\label{eq:finite-factor-error}
 E_i=\frac{\sqrt r(\tau+\epsilon)}{\norm{b_i}}\norm{Y_i}
    +\frac{\sqrt r(\tau+\epsilon)}{\norm{a_i}}\norm{X_i}
    +\frac{\epsilon\norm{X_i}\norm{Y_i}}{\norm{a_i}\norm{b_i}}
\end{equation}
when $a_i,b_i\ne0$; otherwise define $E_i=\norm{X_i}\norm{Y_i}$.

\begin{theorem}[Finite samples with controlled factors]\label{thm:finite-approximation}
If $4r\epsilon+\sum_iE_i<h^2$, then
\[
 N-1\leq\sum_i(m_i-2)_+.
\]
In particular, if all factors of dimension at least three have dimension
at most an integer $d\geq3$, at least
$\lceil(N-1)/(d-2)\rceil$ such factors are required.
\end{theorem}
\begin{proof}
Fix a block with $a,b\ne0$. Nonnegativity of every summand gives
$0\leq\ip a b\leq\epsilon$ and
\[
 \norm{X^*b},\norm{Y^*a}\leq\sqrt r(\tau+\epsilon).
\]
Indeed, each coordinate is the difference of two numbers in
$[0,\tau+\epsilon]$. Set $X_0=P_{b^\perp}X$ and
$Y_0=P_{a^\perp}Y$. Then
\begin{equation}\label{eq:projection-two-term}
 \norm{X^*Y-X_0^*Y_0}
 \leq\frac{\norm{X^*b}}{\norm b}\norm Y
       +\norm X\frac{\norm{Y^*a}}{\norm a}.
\end{equation}
For $m\geq2$, the product $P_{b^\perp}P_{a^\perp}$ has singular
values $1$ with multiplicity $m-2$, together with
$\gamma=\ip a b/(\norm a\norm b)$ and $0$. This follows by
splitting off $(\spanop\{a,b\})^\perp$ and computing the remaining
plane; the parallel case gives the same list with $\gamma=1$.
Deleting the singular direction of size $\gamma$ therefore approximates
$X_0^*Y_0$ by a rank at most $m-2$ matrix with error at most
$\gamma\norm X\norm Y$. Equation~\eqref{eq:projection-two-term}
proves the bound $E_i$. For $m=1$, $X_0=Y_0=0$, so the first two
terms of $E_i$ suffice. If a base factor is zero, the zero approximant
has error at most $\norm{X_i}\norm{Y_i}=E_i$.

The mixed difference matrix of the kernel
$k(v,w)=\sum_i\ip{A_i(v)}{B_i(w)}$, with entries
$k(v_j,v_\ell)-k(v_j,v_0)-k(v_0,v_\ell)+k(v_0,v_0)$, is
$\sum_iX_i^*Y_i$.
For the exact sphere slack it is
$D_*=-h^2I_r-\tau^2\mathbf1\mathbf1^*$, whose smallest singular value
is at least $h^2$. Each entry combines four errors from
\eqref{eq:finite-slack-error}, so the operator error is at most
$4r\epsilon$.
The sum of the block approximants has rank at most $\sum_i(m_i-2)_+$
and lies within $4r\epsilon+\sum_iE_i$ of $D_*$. A singular matrix
cannot lie at operator distance less than $h^2$ from $D_*$: apply it
to a unit vector in its kernel. This proves the assertion.
\end{proof}

A simpler sufficient condition is available when there are $k$ blocks,
$\norm{a_i},\norm{b_i}\geq\rho>0$, and
$\norm{X_i},\norm{Y_i}\leq Lh$. If $\epsilon\leq h^2$, then
$\tau\leq h^2$, and it suffices that
\begin{equation}\label{eq:finite-uniform-threshold}
 \frac{4r\epsilon}{h^2}
 +\frac{4k\sqrt rLh}{\rho}
 +\frac{k\epsilon L^2}{\rho^2}<1.
\end{equation}
For fixed $r,k,L,\rho$, the choice $h$ of order $\epsilon^{1/3}$
makes the left side tend to zero, provided the bounds on
$\norm{X_i},\norm{Y_i}$ hold for the shrinking samples.

If $(1-\eta)\B_2^N\subset C\subset(1+\eta)\B_2^N$, $0<\eta<1$,
choose radial boundary points $x(v)=\rho_C(v)v$ and
$y(w)=\rho_{C^\circ}(w)w$. Their radial lengths lie in
$[1-\eta,1+\eta]$ and $[(1+\eta)^{-1},(1-\eta)^{-1}]$.
Consequently
\[
 |\,1-\ip{x(v)}{y(w)}-(1-\ip vw)\,|
 \leq\frac{2\eta}{1-\eta}.
\]
The reduction and certificate argument of Lemma~\ref{lem:reduction}
provides a factorization of this slack from any exact affine lift of $C$,
including at radial points that are not extreme. Thus metric accuracy
gives~\eqref{eq:finite-slack-error} with $\epsilon=2\eta/(1-\eta)$,
but not the factor bounds in \eqref{eq:finite-factor-error}
or~\eqref{eq:finite-uniform-threshold}.

\subsection{Smooth near-complementarity}
For $r_0>0$, $\delta,H\geq0$, define
\begin{equation}\label{eq:psi-interpolation}
 \psi_{r_0}(\delta,H)=\inf_{0<s\leq r_0}
             \left(\frac{\delta}{s}+\frac{Hs}{2}\right).
\end{equation}
Because the radius is finite, this covers $H=0$, where the value is
$\delta/r_0$. If $H>0$ and $\sqrt{2\delta/H}\leq r_0$, it equals
$\sqrt{2H\delta}$; for $\delta=0$ the value is zero by a limit.

\begin{lemma}[Quantitative near-complementarity]\label{lem:near-complementarity}
Let $A:U\to K$, $B:U\to K^*$ be $C^2$, where $K\subset E$ is a
proper cone of dimension $m$ and $U$ contains the closed Euclidean ball
of radius $r_0$ about zero. Put
$a=A(0)$, $b=B(0)$, $X=DA(0)$, $Y=DB(0)$, and
$\delta=\ip a b$. Suppose $a,b\ne0$ and let
\[
 H_A=\sup_{\norm u=1,\,|t|\leq r_0}
       |\ip{D^2A(tu)[u,u]}b|,\qquad
 H_B=\sup_{\norm u=1,\,|t|\leq r_0}
       |\ip a{D^2B(tu)[u,u]}|.
\]
Write $U_1=\norm X$, $V_1=\norm Y$ and
\[
 \alpha=\frac{\psi_{r_0}(\delta,H_A)}{\norm b},\qquad
 \beta=\frac{\psi_{r_0}(\delta,H_B)}{\norm a},\qquad
 \gamma=\frac{\delta}{\norm a\norm b}.
\]
For $m\geq2$, $X^*Y$ is within
\begin{equation}\label{eq:smooth-factor-error}
 e=U_1\beta+V_1\alpha+U_1V_1\gamma
\end{equation}
of a matrix of rank at most $m-2$. For a ray, the zero approximant
has error at most $e=\alpha\beta$. If either base factor is zero,
$X^*Y=0$.
\end{lemma}
\begin{proof}
For a unit $u$, the function $f(t)=\ip{A(tu)}b$ is nonnegative,
$f(0)=\delta$, and $|f''|\leq H_A$. Taylor's inequality at $-s$ or
$s$, with sign opposite to $f'(0)$, gives
$|f'(0)|\leq\delta/s+H_As/2$. Hence
$\norm{X^*b}\leq\psi_{r_0}(\delta,H_A)$, and similarly for $Y^*a$.
Apply~\eqref{eq:projection-two-term} and its singular-value truncation
to obtain~\eqref{eq:smooth-factor-error}. For $m=1$ the normal
components are the entire derivatives, giving $\alpha\beta$.
If $a=0$, every two-sided derivative of $A$ lies in
$K\cap(-K)=\{0\}$, so $X=0$; the dual case is identical.
\end{proof}

The bound~\eqref{eq:smooth-factor-error} implies the more conservative
bound $U_1\beta+V_1\alpha+\alpha\beta\gamma+U_1V_1\gamma$, obtained by
expanding both normal components separately. Keeping one
projected factor in~\eqref{eq:projection-two-term} removes the extra
term. Both bounds are unchanged when $A$ is multiplied by a positive
constant and $B$ by its reciprocal.

\begin{theorem}[Smooth approximation with controlled factors]\label{thm:smooth-approximation}
Let $\B_2^N\subset C\subset(1+\epsilon)\B_2^N$ and suppose a local
$C^2$ factorization on a sphere chart $\sigma:U\to\Sph^{N-1}$ satisfies
\begin{equation}\label{eq:approx-support-slack}
 \sum_i\ip{A_i(u)}{B_i(v)}
 =h_C(\sigma(v))-\ip{\sigma(u)}{\sigma(v)},\qquad
 D\sigma(0)^*D\sigma(0)=I_{N-1}.
\end{equation}
Apply Lemma~\ref{lem:near-complementarity} in each factor span, of
dimension $s_i$, with resulting error $e_i$. If $\sum_i e_i<1$, then
\[
 N-1\leq\sum_i(s_i-2)_+.
\]
For a definable exact lift of $C$, after minimal-face reduction, such
local $C^2$ factorizations exist on some nonempty open subset of the
sphere. The condition $\sum_i e_i<1$ is an additional hypothesis.
\end{theorem}
\begin{proof}
At the chart center, nonnegativity gives $\delta_i\geq0$ and
$\sum_i\delta_i=h_C(\sigma(0))-1\leq\epsilon$.
Mixed differentiation of~\eqref{eq:approx-support-slack} gives the exact
identity $-\sum_iDA_i(0)^*DB_i(0)=I_{N-1}$, because the support term
depends only on $v$. Lemma~\ref{lem:near-complementarity} approximates
this identity by a matrix of rank at most $\sum_i(s_i-2)_+$ with
error at most $\sum_i e_i$. Error less than one makes that matrix
invertible.

For the existence statement, write the reduced lift as
$C=\pi\{z\in K:Mz=b_0\}$ with relative Slater feasibility, absorbing
an affine output offset into the dual objective if necessary. Choose
the minimum-norm feasible lift of $\sigma(u)$ and the minimum-norm
optimal dual multiplier $\lambda(v)$ for maximizing
$\ip{\sigma(v)}{\pi z}$. Both choices exist in closed nonempty fibers;
Slater dual attainment supplies the second fiber. With
$B(v)=M^*\lambda(v)-\pi^*\sigma(v)$, weak duality and attainment
give~\eqref{eq:approx-support-slack}. Definable choice by these unique
minimizers and $C^2$ cell decomposition give a common open patch.
We do not claim a uniform patch size or uniform derivative bounds
across families of lifts.
\end{proof}

For example, suppose there are at most $k$ nonzero selected blocks,
$\norm{a_i},\norm{b_i}\geq\mu>0$,
$\norm{DA_i},\norm{DB_i}\leq L$, and $H_{A,i},H_{B,i}\leq H$,
where $H>0$ and $\sqrt{2\epsilon/H}\leq r_0$. The sum of the errors
for blocks of dimension at least two is at most
\begin{equation}\label{eq:uniform-smooth-error}
 \mathcal E_k(\epsilon)
 =\frac{2L\sqrt{2Hk\epsilon}}{\mu}
       +\frac{L^2\epsilon}{\mu^2}.
\end{equation}
This follows from $\sum_i\sqrt{\delta_i}\leq\sqrt{k\epsilon}$.
Rays add at most $2H\epsilon/\mu^2$. Zero selected factors add zero.
Thus a bounded number of blocks, factor norms bounded away from zero,
bounded derivatives, and a fixed chart radius give error
$O(\sqrt\epsilon)$. Conversely, a family with
$\sum_i(s_i-2)_+<N-1$ must have
$\mathcal E_k(\epsilon)+2H\epsilon/\mu^2\geq1$ (omit the last term
when no rays occur). If the terms other than the square-root term sum
to at most $1/2$, then
\[
 L\sqrt H\geq\frac{\mu}{4\sqrt{2k\epsilon}}.
\]
This concerns the selected factors, not an intrinsic condition number
of a Newton system. Global consequences under stronger selection
hypotheses appear in Appendix~\ref{app:approximation-refinements}.

\subsection{Why approximating the convex body is not enough}
For every $\epsilon>0$, choose a finite $\delta$-net
$\mathcal U\subset\Sph^{N-1}$ with
$(1-\delta^2/2)^{-1}\leq1+\epsilon$. The polytope
\[
 P=\{x:\ip ux\leq1\text{ for all }u\in\mathcal U\}
\]
contains $\B_2^N$. For $x\ne0$ choose $u$ within $\delta$ of
$x/\norm x$. Then $\ip u{x/\norm x}\geq1-\delta^2/2$, so
$\norm x\leq1+\epsilon$ for $x\in P$. One nonnegative slack per
inequality gives an exact lift of $P$ by rays only. Therefore no
bound on the Hausdorff or multiplicative error alone, however small,
forces $N-1\leq\sum_i(m_i-2)_+$. The number of scalar inequalities is a separate
resource. For lifted polyhedral approximations of the Lorentz cone,
Ben-Tal and Nemirovski prove matching bounds of order
$N\log(1/\epsilon)$ on this number for small $\epsilon$
\cite[Theorem~1.1 and Proposition~3.1]{BTN2001}. If such an
approximation is realized with an orthant of that dimension, the optimal
ambient barrier parameter equals the number of slacks.
This is a property of that realization, not of every cone that lifts
the same polytope.

At a fixed error, the bound can fail even for lifts without rays. For $d\geq3$,
\[
 \B_2^{2(d-1)}\subset
 \B_2^{d-1}\times\B_2^{d-1}
 \subset\sqrt2\B_2^{2(d-1)}.
\]
The outer ratio and Hausdorff distance $\sqrt2-1$ are attained by
pairs of unit vectors. Two $Q_d$ factors represent the product, whereas
the exact ball requires at least
$\lceil(2d-3)/(d-2)\rceil=3$ factors of dimension at most $d$.

Bounded factors and an arbitrarily small sampled slack error do not
replace the hypothesis of Theorem~\ref{thm:finite-approximation}, which
controls the increments relative to the base factors.
For its two-point sample ($r=1$), fix $\alpha>0$ and put
$\beta=\tau/\alpha$. In $\R_+^2$ take
\[
 A(v_0)=B(v_0)=(\alpha,\alpha),\quad
 A(v_1)=(0,\alpha+\beta),\quad B(v_1)=(\alpha+\beta,0).
\]
The kernel values are $2\alpha^2,\tau+\alpha^2,
\tau+\alpha^2,0$, so the error tends to zero. Its mixed difference
is nevertheless $-2\tau=-h^2-\tau^2$, exactly the sphere value,
although $(m-2)_+=0$. Apply the diagonal automorphism
$D=\operatorname{diag}(\sqrt{\beta/\alpha},\sqrt{\alpha/\beta})$
to $A$ and $D^{-1}$ to $B$. All factors become bounded as
$\alpha\downarrow0$ and the pairings stay unchanged, but the
hypothesis still fails. This finite example does not by
itself construct a global fixed-size approximation of a ball.

Approximate slack factorizations and their geometric approximation
properties are classical. Gouveia, Parrilo, and Thomas bound inner and
outer polytope approximations using the maximum column and row Euclidean
norms of a finite slack-matrix error; they also obtain approximate
factorizations from lifted approximations of nested polyhedra
\cite[Sections~2.7 and~4.2]{GPTApprox2015}. Our results add explicit
local conditions on the factors at curved contact points under which the
rank bound $\dim K_i-2$ holds for approximate lifts. We do not claim that
every metric approximation admits such factors, or that violating the
conditions imposes an algorithmic cost.
